% Part: turing-machines
% Chapter: undecidability
% Section: trakhtenbrot

\documentclass[../../../include/open-logic-section]{subfiles}

\begin{document}

\olfileid[fr]{tur}{und}{tra}
\olsection{Le théorème de Trakhtenbrot}

\begin{explain}
Dans la \olref[rep]{sec}, nous avons défini les !!{sentence}s
$!T(M,w)$ et~$!E(M,w)$ associés à une machine de Turing~$M$ et
à une entrée~$w$. Nous avons ensuite établi, dans
le \olref[ver]{lem:valid-if-halt} et
le \olref[ver]{lem:halt-if-valid}, que
$!T(M,w)\lif !E(M,w)$ est valide si et seulement si~$M$
finit par s'arrêter sur~$w$. L'insolubilité du problème de l'arrêt
entraîne donc l'indécidabilité de la validité et de la
satisfaisabilité des !!{sentence}s du premier ordre
(\cref{tur:und:uns:thm:decision-prob,tur:und:uns:cor:undecidable-sat}).

Ces deux notions portent toutefois sur des !!{structure}s
arbitraires, finies ou infinies. On pourrait penser qu'il est plus
facile de décider si !!a{sentence} est satisfaisable dans une
!!{structure} finie, ou valide dans toutes les !!{structure}s
finies. En adaptant la réduction précédente, nous allons montrer
que ce n'est pas le cas.

Revenons à la preuve du \olref[ver]{lem:halt-if-valid}.
Nous y avons construit un modèle~$\Struct M$ de~$!T(M,w)$
décrivant exactement l'exécution de~$M$ sur~$w$. Son domaine,
$\Nat$, était infini. Si~$M$ s'arrête après $k$ étapes, nous
pouvons construire de la même manière un modèle fini
$\Struct{M'}$. Prenons
$\Domain{M'}=\{0,\dots,N\}$ avec
$N=\len{w}+k+2$.\footnote{Correction éditoriale : la borne source $\max(k,\len{w})$ peut confondre le dernier instant ou une case visitée avec le bord où le successeur boucle. Nous plaçons ce bord strictement au-delà de l'entrée, de toutes les positions visitées et de l'arrêt ; la même correction est utilisée pour $!T'(M,w)$.} Interprétons $\prime$ par la fonction successeur
jusqu'à~$N$, puis par une boucle en~$N$ :
\[
\Assign{\prime}{M'}(x)=
\begin{cases}
  x+1 &\text{si $x<N$,}\\
  N &\text{si $x=N$.}
\end{cases}
\]
Interprétons~$<$ par l'ordre usuel auquel nous ajoutons
$N<N$. Les prédicats d'état et de ruban décrivent la véritable
exécution jusqu'à l'arrêt; ils sont faux aux instants ultérieurs.
La borne~$N$ est strictement supérieure au dernier instant~$k$,
à la dernière case de l'entrée et à toute case que la tête a pu
visiter avant l'arrêt. La boucle en~$N$ satisfait l'axiome
$\lforall[x][x<x']$ au bord du domaine, mais aucune transition
réelle n'atteint ce bord. Ainsi, comme dans la preuve du
\olref[ver]{lem:halt-if-valid},
$\Sat{M'}{!T(M,w)}$ et, puisque~$M$ s'arrête sur~$w$,
$\Sat{M'}{!E(M,w)}$. La structure~$\Struct{M'}$ est donc un
modèle fini de $!T(M,w)\land !E(M,w)$ : l'implication a été
remplacée ici par une conjonction.

Nous avons établi une moitié du résultat : si~$M$ s'arrête sur~$w$,
alors $!T(M,w)\land !E(M,w)$ possède un modèle fini. La réciproque
est fausse. Considérons, par exemple, la machine~$M$ dont l'unique
état est~$q_0$ et dont l'unique instruction est
$\delta(q_0,\TMblank)=\tuple{q_0,\TMblank,\TMstay}$.
Sur l'entrée vide~$w=\emptyseq$, elle ne s'arrête jamais : elle
reste dans une boucle sans changer le ruban ni déplacer sa tête.
Toutes les configurations sont identiques. Pourtant, on peut
construire une !!{structure} finie~$\Struct{M''}$ qui satisfait
$!T(M,\emptyseq)\land !E(M,\emptyseq)$ (exercice).
Nous allons modifier~$!T(M,w)$ pour exclure ce type de
!!{structure}s.

\begin{prob}
Soit~$M$ une machine de Turing ayant pour seul état~$q_0$
et pour seule instruction
$\delta(q_0,\TMblank)=\tuple{q_0,\TMblank,\TMstay}$.
Prenez $\Domain{M''}=\{0,1,2\}$,
$\Assign{\prime}{M''}(0)=\Assign{\prime}{M''}(1)=1$,
$\Assign{\prime}{M''}(2)=2$, et\footnote{Corrections éditoriales : l'état cible de l'unique instruction est $q_0$, comme dans l'exemple précédent, et les deux affectations du successeur concernent $\Struct{M''}$ ; la source y écrit respectivement $q$ et $\Struct{M'}$.}
$\Assign{<}{M''}=\{\tuple{0,1},\tuple{1,1},\tuple{2,2}\}$.
Définissez les interprétations de~$\Obj Q_{q_0}$,
$\Obj S_{\TMblank}$ et~$\Obj S_{\TMendtape}$ de façon à rendre
vrais $!T(M,\emptyseq)$ et~$!E(M,\emptyseq)$; expliquez pourquoi
ils le sont. Indication : $\delta(q_0,\TMendtape)$ n'est pas
définie. Faites en sorte que
\begin{align*}
& \Obj Q_{q_0}(\num{1},\num{n})\land
  \Obj S_{\TMendtape}(\num{0},\num{n})\land
  \lforall[x][(\num{0}<x\lif
  \Obj S_\TMblank(x,\num{n}))]
  \text{\quad pour tout $n\in\Nat$}\\
& \lexists[y][(\Obj Q_{q_0}(\num{0},y)\land
  \Obj S_{\TMendtape}(\num{0},y))]
\end{align*}
soient tous deux vrais dans~$\Struct{M''}$.
\end{prob}
\end{explain}

Considérons les !!{sentence}s décrivant le fonctionnement de~$M$
sur l'entrée $w=\sigma_{i_1}\dots\sigma_{i_k}$ :
\begin{enumerate}
\item Les axiomes portant sur les nombres et sur~$<$, comme dans
  la définition de~$!T(M,w)$ (\olref[rep]{sec}).
\item Les axiomes décrivant la configuration initiale, comme
  dans la définition de~$!T(M,w)$.
\item Les axiomes décrivant le passage d'une configuration à la
  suivante.

Gardons la définition précédente de~$!A(x,y)$ et posons
\[
!B(y)\ident\lforall[x][(x<y\lif\eq/[x][y])].
\]
Ici $\eq/[x][y]$ signifie que $x$ est \emph{distinct} de~$y$.
\begin{enumerate}
\item \ollabel{rep-right} Pour chaque instruction
  $\delta(q_i,\sigma)=\tuple{q_j,\sigma',\TMright}$,
  l'!!{sentence} :
\begin{align*}
& \lforall[x][\lforall[y][(
  (\Obj Q_{q_i}(x,y)\land\Obj S_\sigma(x,y))\lif {}]]\\
&\qquad (\Obj Q_{q_j}(x',y')\land
  \Obj S_{\sigma'}(x,y')\land !A(x,y)\land !B(y')))
\end{align*}
\item \ollabel{rep-left} Pour chaque instruction
  $\delta(q_i,\sigma)=\tuple{q_j,\sigma',\TMleft}$,
  l'!!{sentence} :
\begin{align*}
& \lforall[x][\lforall[y][(
  (\Obj Q_{q_i}(x',y)\land\Obj S_\sigma(x',y))\lif {}]]\\
&\qquad (\Obj Q_{q_j}(x,y')\land
  \Obj S_{\sigma'}(x',y')\land
  !A(x',y)\land !B(y')))\land {}\\
& \lforall[y][(
  (\Obj Q_{q_i}(\num{0},y)\land
  \Obj S_\sigma(\num{0},y))\lif {}]\\
&\qquad (\Obj Q_{q_j}(\num{0},y')\land
  \Obj S_{\sigma'}(\num{0},y')\land
  !A(\num{0},y)\land !B(y')))
\end{align*}
\item \ollabel{rep-stay} Pour chaque instruction
  $\delta(q_i,\sigma)=\tuple{q_j,\sigma',\TMstay}$,
  l'!!{sentence} :
\begin{align*}
& \lforall[x][\lforall[y][(
  (\Obj Q_{q_i}(x,y)\land\Obj S_\sigma(x,y))\lif {}]]\\
&\qquad (\Obj Q_{q_j}(x,y')\land
  \Obj S_{\sigma'}(x,y')\land !A(x,y)\land !B(y')))
\end{align*}
\end{enumerate}
Ces !!{sentence}s reprennent les clauses de transition de~$!T(M,w)$
en leur ajoutant $!B(y')$.\footnote{Corrections éditoriales : dans le premier cas du déplacement gauche, la source emploie le terme de conservation $!A(x,y)$ et omet $!B(y')$. Nous écrivons $!A(x',y)$, qui exclut la case écrite, et ajoutons la condition de fraîcheur à ce cas aussi.} Cette dernière formule assure que
l'instant~$y'$ de la nouvelle configuration diffère de tous les
instants précédents~$\Obj 0$, $\Obj 0'$, \dots{}.
% Conditions de cohérence proposées dans la source, mais inactives :
% \item Énoncés assurant des conditions de cohérence :
% \begin{enumerate}
%   \item\ollabel{rep-con-symb}%
%   !!^a{sentence} affirmant qu'aucune case ne porte deux symboles
%   distincts au même instant :
%   \[\lforall[x][\lforall[y][(\Obj S_\sigma(x,y) \lif
%   \lnot\Obj S_{\sigma'}(x,y))]],\]
%   pour deux symboles distincts $\sigma,\sigma'$.
%   \item\ollabel{rep-con-state}%
%   !!^a{sentence} affirmant que la machine n'a pas deux états
%   distincts au même instant :
%   \[\lforall[x][\lforall[y][(\Obj Q_q(x,y)\lif
%   \lforall[z][\lnot\Obj Q_{q'}(z,y))]],\]
%   pour deux états distincts $q,q'$.
%   \item\ollabel{rep-con-tape}%
%   !!^a{sentence} affirmant que la tête ne lit pas deux cases
%   distinctes au même instant :
%   \[\lforall[x][\lforall[y][(\Obj Q_q(x,y)\lif
%   \lforall[z][\eq/[z][y]\lif\lnot\Obj Q_{q'}(z,y))]],\]
%   pour deux états $q,q'$.
% \end{enumerate}
\end{enumerate}
Notons $!T'(M,w)$ la conjonction de tous les !!{sentence}s
précédents, pour la machine~$M$ et l'entrée~$w$.

\begin{lem}\ollabel{lem:halts-sat}
Si~$M$ s'arrête sur l'entrée~$w$, alors
$!T'(M,w)\land !E(M,w)$ possède un modèle fini.
\end{lem}

\begin{proof}
Supposons que~$M$ s'arrête pour la première fois après $k$ étapes.
Construisons $\Struct{M'}$ comme la structure de la preuve du
\olref[ver]{lem:halt-if-valid}, sauf que son domaine est fini :
\begin{align*}
\Domain{M'}&=\{0,\dots,N\},\\
\Assign{\prime}{M'}(x)&=
\begin{cases}
  x+1 &\text{si $x<N$,}\\
  N &\text{si $x=N$,}
\end{cases}\\
\tuple{x,y}\in\Assign{<}{M'}
&\text{ si et seulement si $x<y$ ou $x=y=N$,}
\end{align*}
où $N=\len{w}+k+2$. Les autres symboles sont interprétés
comme dans la structure décrivant l'exécution réelle de~$M$
jusqu'à son arrêt; au-delà, les relations d'état et de ruban
sont vides.

Vérifions les points propres au domaine fini. L'axiome
$\lforall[x][x<x']$ est satisfait, y compris pour $x=N$
grâce à~$N<N$; la relation~$<$ reste transitive.
La case~$N$ n'est jamais lue par la tête et l'instant~$N$
n'est pas atteint avant l'arrêt. Ainsi, la boucle du
successeur n'altère aucune transition réelle ni la conservation
des symboles dans~$!A$. Pour chaque transition exécutée à
l'instant~$y<k$, le nouvel instant~$y+1$ est inférieur à~$N$;
la relation~$<$ est l'ordre usuel en deçà de~$N$, de sorte que
$!B(y')$ est vrai. Les axiomes de la configuration initiale et
les autres clauses de transition sont satisfaits comme dans
$\Struct M$. Enfin, la configuration terminale satisfait
$!E(M,w)$. Nous avons donc
$\Sat{M'}{!T'(M,w)\land !E(M,w)}$.
\end{proof}

\begin{prob}
Complétez la preuve du \olref[tur][und][tra]{lem:halts-sat}
en vérifiant directement
$\Sat{M'}{!T'(M,w)\land !E(M,w)}$, notamment la clause
$!B(y')$ pour chaque transition effectivement exécutée.\footnote{Correction éditoriale : l'exercice source demande seulement $!T(M,w)\land !E(M,w)$. Le lemme porte sur $!T'(M,w)\land !E(M,w)$ ; la condition de fraîcheur doit donc aussi être vérifiée.}
\end{prob}

\begin{lem}\ollabel{lem:sat-halts}
Si $!T'(M,w)\land !E(M,w)$ possède un modèle fini, alors~$M$
s'arrête sur l'entrée~$w$.
\end{lem}

\begin{proof}
Raisonnons par contraposition. Supposons que~$M$ ne s'arrête pas
sur~$w$. Montrons qu'aucun modèle de
$!T'(M,w)\land !E(M,w)$ ne peut être fini. Soit~$\Struct M$
un tel modèle, s'il en existe un.

Comme dans la preuve du \olref[ver]{lem:config}, les clauses
initiales et de transition entraînent
$!T'(M,w)\Entails !C(M,w,n)$ pour tout~$n\in\Nat$.
L'exécution ne s'arrêtant jamais, la transition de l'instant
$n-1$ fournit en outre
$!T'(M,w)\Entails !B(\num{n})$ pour chaque~$n\ge1$.
D'après la \olref[rep]{prop:mlessk}, les axiomes d'ordre donnent
$!T'(M,w)\Entails\num{k}<\num{n}$ pour tout $k<n$.
Or $!B(\num{n})$ entraîne
$\num{k}<\num{n}\lif\eq/[\num{k}][\num{n}]$.
Par conséquent,
$\Sat{M}{\eq/[\num{k}][\num{n}]}$ dès que $k<n$ :
les termes $\num{0},\num{1},\dots$ prennent des valeurs
deux à deux distinctes dans~$\Struct M$. Son domaine est donc
infini. Aucun modèle fini de
$!T'(M,w)\land !E(M,w)$ n'existe.
\end{proof}

\begin{prob}
Complétez la preuve du \olref[tur][und][tra]{lem:sat-halts}
en montrant que, si~$M$ ne s'arrête pas sur~$w$,
$!T'(M,w)\Entails !B(\num{n})$ pour tout~$n\ge1$.
Expliquez pourquoi le cas $n=0$ n'est pas nécessaire.\footnote{Correction éditoriale : $!B(\num{n})$ vient de la transition précédente ; la source l'affirme aussi pour $n=0$, alors qu'aucune condition de ce genre n'est imposée avant le premier pas. Pour distinguer $\num{k}$ et $\num{n}$ avec $k<n$, on a toujours $n\ge1$.}
\end{prob}

\begin{thm}[Théorème de Trakhtenbrot]
\ollabel{thm:trakhtenbrodt}
Le problème de savoir si un !!{sentence} arbitraire de logique
du premier ordre possède un modèle fini, autrement dit s'il
est finiment satisfaisable, est indécidable.
\end{thm}

\begin{proof}
Supposons qu'une machine de Turing~$F$ décide la satisfaisabilité
finie. Étant données une machine de Turing~$M$ et une entrée~$w$,
nous pouvons calculer l'!!{sentence}
$!T'(M,w)\land !E(M,w)$, puis demander à~$F$ s'il possède un
modèle fini. D'après les
\cref{tur:und:tra:lem:halts-sat,tur:und:tra:lem:sat-halts},
la réponse est affirmative si et seulement si~$M$ s'arrête sur~$w$.
La machine~$F$ permettrait donc de résoudre le problème de l'arrêt,
ce qui est impossible.
\end{proof}

\begin{cor}\ollabel{cor:fproof-incomp}%
Il n'existe aucun système de !!{derivation} dont les
!!{derivation}s soient effectivement énumérables et qui soit
à la fois correct et complet pour la \emph{validité finie}.\footnote{Précision éditoriale : l'exclusion concerne les systèmes effectifs. Sans énumérabilité effective des dérivations, l'ensemble de tous les énoncés valides dans les structures finies pourrait lui-même être pris pour ensemble d'axiomes.}
Un tel système vérifierait, pour tout énoncé~$!B$,
$\Proves !B$ si et seulement si
$\Sat{M}{!B}$ pour toute !!{structure} finie~$\Struct M$.
\end{cor}

\begin{proof}
La preuve est laissée en exercice.
\end{proof}

\begin{prob}
Démontrez le \olref[tur][und][tra]{cor:fproof-incomp}.
Remarquez que~$!B$ est vrai dans toute !!{structure} finie
si et seulement si~$\lnot !B$ n'est pas finiment satisfaisable.
Expliquez pourquoi la satisfaisabilité finie est semi-décidable
au sens du \olref[tur][und][uns]{thm:valid-ce}, en énumérant
les interprétations finies pertinentes. Montrez ensuite que,
si les !!{derivation}s d'un système correct et complet pour
la validité finie étaient effectivement énumérables, la
satisfaisabilité finie serait décidable.
\end{prob}

\end{document}
